\subsection{AI for AI: The Second Axis of Self-Improvement}
\label{sec:rsi-vertical}

Every campaign so far produces something outside the runtime: a kernel, a manuscript,
a crystal structure, a chip. This one points inward. The artifact is the machinery
that produces models, which is the machinery that produces \systemname{} itself.

Self-improvement comes in two kinds, and conflating them is why claims in this area
are hard to evaluate. \emph{Parameter-invariant} self-improvement changes what the
system knows, retrieves, and obeys while the weights stay fixed: memories, skills,
procedures, verifiers, routing decisions, and retained failures. This report measures
it directly---Section~\ref{sec:method} specifies the admission gate, and
Section~\ref{sec:results} shows mature waves using \SWEProTokenReduction\% fewer
solve input Tokens and \SWEProTimeReduction\% less active time per task than startup,
under a backbone that never changed. Systems in this family evolve their
context, their retained experience, or their own scaffolding code under empirical
feedback~\citep{zhang2025ace,yin2025godelagent,zhang2025darwingodel}. \emph{Parameter-changing}
self-improvement alters the weights themselves. The two are complementary rather than competing: the
first compounds within a fixed capability ceiling, and the second is what moves the
ceiling.

\systemname{} operates on the first axis today, and that is the axis where evidence
is actually obtainable, because a state update has an owner, a precondition, and an
audit trail. But the ceiling is real. No amount of accumulated skill makes a frozen
backbone reason better than it can, and a system confined to that axis improves
asymptotically toward its own limit. Reaching the second axis requires participating
in the production of weights, which means owning a pretraining stack: data pipeline,
training loop, stability, throughput, failure recovery, and the judgement to decide
which of those is the current bottleneck.

\systemname{} built and ran that stack. On an eight-GPU B200 node, the runtime took a
1B-parameter model from scratch through data-quality and pipeline construction,
training stability and speed, GPU utilization and bottleneck analysis, engineering
quality and failure recovery, and successive optimization rounds---with the operator
driving it having prior experience only in post-training, not in pretraining
infrastructure. The pipeline runs end to end and continues under the runtime's own
supervision.

\paragraph{What this is evidence for, and what it is not.}
The claim here is deliberately narrow, because the interesting failure would be to
overstate it. Building a pretraining stack is not yet parameter-changing
self-improvement; it is the entry condition for it. Until a system can produce
weights, the second axis is unavailable in principle, not merely unattempted.

Three steps remain, and each is genuinely hard. The trained model is not yet good
enough to serve as a role backbone. The training data is a public corpus rather than
the runtime's own retained trajectories, which
Section~\ref{sec:runtime-to-training} argues is the natural source and which the
first axis has been accumulating all along. And no model produced this way has been
placed into a role and measured for whether the runtime that built it improved. Only
that last measurement would demonstrate a closed loop rather than a capability.

The two axes are not independent, which is the point of reporting them together. The
trajectories that parameter-invariant self-evolution retains---objectives, actions,
measurements, revisions, verdicts, and the routes that were refuted---are exactly the
supervision a training stage would consume. A runtime that improves its own state
while producing a structured record of how it did so is, in effect, generating its
own training corpus as a side effect of operating. Owning the stack that could
consume it is what makes that observation more than rhetorical.

\paragraph{Why this vertical is a hard test of unattended operation.}
Pretraining is the least forgiving setting in this report for a Driver that cannot
judge its own progress. A multi-day run commits resources irreversibly, and its
characteristic failures---a data pipeline defect, a silent throughput regression, an
instability that only manifests after a loss curve has looked healthy for
hours---surface long after the decision that caused them. There is no fast verifier
to consult, no test that turns green, and the cost of noticing late is measured in
GPU-days. This is precisely the regime in which Question~Q1 of
Section~\ref{sec:introduction}---is this actually working---cannot be answered by
checking an exit code, and where a runtime that persists without judging is worse
than one that stops.
